% Generated by make_examples.py from the model. Do not edit by hand.
\subsection{Inputs}
Table~\ref{tab:in-Cirrus_SR22T} lists the aircraft data and the day's conditions. The example uses Wichita Dwight D. Eisenhower National (KICT) runway 01L.
\begin{table}[H]\centering\small
\begin{tabular}{@{}ll@{}}\toprule \textbf{Input} & \textbf{Value}\\\midrule
Maximum takeoff weight (MTOW) & \SI{3600}{lb} = \SI{1632.9}{kg}\\
Maximum landing weight (MLW) & \SI{3600}{lb} = \SI{1632.9}{kg}\\
Reference wing area, $S$ & \SI{13.46}{m^2} (\SI{144.9}{ft^2})\\
$C_{D_0}$ / $k$ & \num{0.027} / \num{0.045}\\
$C_{L,\max}$ takeoff / landing & \num{1.60} / \num{2.04}\\
$C_{L,\text{gnd}}$ & \num{0.40}\\
$\mu_\text{roll}$ / $\mu_\text{brake}$ & \num{0.04} / \num{0.40}\\
Sea-level static thrust, $T_\text{SLS}$ & \SI{3580}{N} (\SI{805}{lbf})\\
Number of engines, $N_e$ & 1\\
Maximum power, $P_\text{max}$ & \SI{234.9}{kW} (\SI{315}{hp})\\
Propeller efficiency, $\eta_p$ & \num{0.72}\\
Turbocharger critical altitude & \SI{25000}{ft}\\
Field elevation, $h$ & \SI{1333}{ft} = \SI{406.3}{m}\\
Outside air temperature (OAT) & \SI{30}{degreeCelsius}\\
Reported wind & \ang{20} / \SI{10}{kt}\\
Runway true heading & \ang{20}\\
Runway slope (direction of travel) & \num{0.165}\,\%\\
Runway condition & Dry\\
TODA = ASDA / LDA & \SI{10302}{ft} / \SI{10302}{ft}\\
\bottomrule\end{tabular}\caption{Cirrus SR22T: example inputs.}\label{tab:in-Cirrus_SR22T}\end{table}
\subsection{Step 1: Atmosphere}
Standard temperature at the field elevation, then the ISA deviation from the OAT (\eqref{eq:isa-temp} and \eqref{eq:isa-dev}):
\begin{align*}
\Theta_\text{std} &= \num{288.15} - \num{0.0065}\times\num{406.3} = \SI{285.51}{K}\\
\Delta T &= \num{30.0} - (\num{285.51} - \num{273.15}) = \SI{17.64}{degreeCelsius}\quad\text{(ISA }+17.6\text{)}
\end{align*}
Pressure, temperature, density, density ratio and speed of sound (\eqref{eq:isa-p}--\eqref{eq:isa-a}):
\begin{align*}
p &= \num{101325}\left(\frac{\num{285.51}}{\num{288.15}}\right)^{\num{5.2559}} = \SI{96538}{Pa}\\
\Theta &= \num{285.51} + \num{17.64} = \SI{303.15}{K}\\
\rho &= \frac{\num{96538}}{\num{287.05}\times\num{303.15}} = \SI{1.1094}{kg/m^3}\\
\sigma &= \frac{\num{1.1094}}{\num{1.225}} = \num{0.9056}\\
a &= \sqrt{\num{1.4}\times\num{287.05}\times\num{303.15}} = \SI{349.0}{m/s}
\end{align*}
\subsection{Step 2: Wind}
Runway components (\eqref{eq:wind-hw}--\eqref{eq:wind-xw}), then the regulatory factor \eqref{eq:wind-perf}:
\begin{align*}
V_w &= \num{10}\cos(\ang{20} - \ang{20}) = \SI{10.0}{kt}\qquad V_x = \num{10}\sin(\ang{0}) = \SI{0.0}{kt}\\
V_{w,\text{perf}} &= \num{0.5}\times\num{10.0} = \SI{5.0}{kt} = \SI{2.57}{m/s}
\end{align*}
All following steps use the performance headwind of \SI{5.0}{kt}.
\subsection{Step 3: Weight and Takeoff Speeds}
\begin{align*}
W &= m\,g_0 = \num{1632.9}\times\num{9.80665} = \SI{16014}{N}\\
V_{S1} &= \sqrt{\frac{2\times\num{16014}}{\num{1.1094}\times\num{13.46}\times\num{1.60}}} = \SI{36.61}{m/s} = \SI{71.2}{kt}\\
V_R &= \num{1.10}\,V_{S1} = \SI{40.27}{m/s} = \SI{78.3}{kt}\\
V_2 &= \num{1.20}\,V_{S1} = \SI{43.93}{m/s} = \SI{85.4}{kt}\\
V_\text{TR} &= \num{1.15}\,V_{S1} = \SI{42.10}{m/s} = \SI{81.8}{kt}
\end{align*}
\subsection{Step 4: Forces on the Ground Roll}
Thrust \eqref{eq:T-prop}, turbocharged below its critical altitude so the power is not reduced ($\sigma_P = 1$):
\begin{equation*}
T(V) = \min\left(\num{1.0000}\times\num{3580},\ \frac{\num{0.72}\times\num{1.0000}\times\num{234895}}{V}\right)
\end{equation*}
\Needspace{16\baselineskip}
Forces at brake release and at $V_R$ (\eqref{eq:eom-accel}, $C_L = C_{L,\text{gnd}} = \num{0.40}$, slope force $W\sin\theta$):
\begin{table}[H]\centering\small
\begin{tabular}{@{}lrr@{}}\toprule \textbf{Quantity} & \textbf{At $V_\text{TAS} = V_w$ (start)} & \textbf{At $V_R$}\\\midrule
Dynamic pressure $q$ (Pa) & \num{3.7} & \num{899.6}\\
Lift $L$ (N) & \num{20} & \num{4844}\\
Drag $D$ (N) & \num{2} & \num{414}\\
Thrust $T$ (N) & \num{3580} & \num{3580}\\
Rolling friction $\mu_\text{roll}(W-L)$ (N) & \num{640} & \num{447}\\
Slope force $W\sin\theta$ (N) & \num{26} & \num{26}\\
Acceleration $a$ (\si{m/s^2}) & \num{1.783} & \num{1.649}\\
\bottomrule\end{tabular}\caption{Cirrus SR22T: forces on the ground roll.}\end{table}
\subsection{Step 5: Ground Roll by Numerical Integration}
Equation~\eqref{eq:eom-accel} is integrated with RK45 from rest until $V_\text{TAS} = V_R$. Selected points from the solver:
\begin{table}[H]\centering\small
\begin{tabular}{@{}rrrr@{}}\toprule \textbf{$t$ (s)} & \textbf{$V_\text{TAS}$ (kt)} & \textbf{$a$ (\si{m/s^2})} & \textbf{$s$ (ft)}\\\midrule
\num{0.0} & \num{5.0} & \num{1.783} & \num{0}\\
\num{2.1} & \num{12.3} & \num{1.781} & \num{13}\\
\num{6.1} & \num{26.1} & \num{1.769} & \num{109}\\
\num{10.1} & \num{39.8} & \num{1.749} & \num{298}\\
\num{14.1} & \num{53.3} & \num{1.721} & \num{579}\\
\num{18.1} & \num{66.5} & \num{1.686} & \num{949}\\
\midrule \num{21.73} & \num{78.3} ($V_R$) & \num{1.649} & \num{1361}\\
\bottomrule\end{tabular}\caption{Cirrus SR22T: ground-roll time history (RK45 output).}\end{table}
Rotation for $t_\text{rot}$ \eqref{eq:s-rot}:
\begin{align*}
s_\text{rot} &= (\num{40.27} - \num{2.57})\times\num{1} = \SI{37.7}{m}\\
s_\text{ground} &= \num{414.9} + \num{37.7} = \SI{452.6}{m} = \SI{1485}{ft}
\end{align*}
\subsection{Step 6: Air Distance to the Screen Height}
Climb at $V_2$ with lift equal to weight (\eqref{eq:cl-climb}--\eqref{eq:sin-gamma}); screen height \SI{50}{ft} = \SI{15.24}{m}:
\begin{align*}
q_{V_2} &= \tfrac12\times\num{1.1094}\times\num{43.93}^2 = \SI{1070.6}{Pa}, \qquad C_{L,\text{climb}} = \frac{\num{16014}}{\num{1070.6}\times\num{13.46}} = \num{1.111}\\
D &= \num{1070.6}\times\num{13.46}\times(\num{0.027} + \num{0.045}\times\num{1.111}^2) = \SI{1190}{N}, \qquad T = \SI{3580}{N}\\
\sin\gamma &= \frac{\num{3580} - \num{1190}}{\num{16014}} = \num{0.1493}, \qquad \gamma = \SI{0.1498}{rad}, \qquad \text{gradient} = \num{15.1}\,\%
\end{align*}
Transition arc (\eqref{eq:R-tr}--\eqref{eq:h-tr}):
\begin{align*}
R_\text{TR} &= \frac{\num{42.10}^2}{\num{9.807}\times(\num{1.2} - 1)} = \SI{903.8}{m}\\
h_\text{TR} &= \num{903.8}\,(1 - \cos \num{0.1498}) = \SI{10.12}{m}
\end{align*}
Since $h_\text{TR} < h_\text{scr}$, the arc is followed by a steady climb \eqref{eq:s-air-climb}:
\begin{equation*}
s_\text{air} = \num{903.8}\sin \num{0.1498} + \frac{\num{15.24} - \num{10.12}}{\tan \num{0.1498}} = \num{134.9} + \num{33.9} = \SI{168.8}{m}
\end{equation*}
Corrected for wind \eqref{eq:s-air-wind}, then the takeoff distance \eqref{eq:tod}:
\begin{align*}
s_\text{air} &= \num{168.8}\times\frac{\num{42.10} - \num{2.57}}{\num{42.10}} = \SI{158.5}{m}\\
\text{TOD}_\text{AEO} &= \num{452.6} + \num{158.5} = \SI{611.1}{m} = \SI{2005}{ft}
\end{align*}
\subsection{Step 7: Engine-Failure Case}
The aircraft has one engine, so there is no accelerate-go case and no balanced field length. The required takeoff distance is the all-engines TOD.
\subsection{Step 8: Landing}
Landing at MLW on the same runway, no reverse thrust credit (\eqref{eq:vref}--\eqref{eq:ldr}):
\begin{align*}
V_{S0} &= \sqrt{\frac{2\times\num{16014}}{\num{1.1094}\times\num{13.46}\times\num{2.04}}} = \SI{32.42}{m/s} = \SI{63.0}{kt}\\
V_\text{REF} &= \num{1.30}\,V_{S0} = \SI{81.9}{kt}, \qquad V_\text{TD} = \num{1.15}\,V_{S0} = \SI{72.5}{kt}\\
s_\text{air} &= \frac{\num{15.24}}{\tan\ang{3}}\times\frac{\num{42.15} - \num{2.57}}{\num{42.15}} = \SI{273.1}{m}\\
s_\text{free} &= (\num{37.29} - \num{2.57})\times\num{2} = \SI{69.4}{m}\\
s_\text{brake} &= \SI{171.9}{m} \quad\text{(\eqref{eq:eom-brake} integrated to rest, \num{9.5}\,s)}\\
\text{LD} &= \num{273.1} + \num{69.4} + \num{171.9} = \SI{514.4}{m} = \SI{1688}{ft}
\end{align*}
\subsection{Step 9: Runway Check}
\begin{table}[H]\centering\small
\begin{tabular}{@{}lrlrrl@{}}\toprule \textbf{Required} & \textbf{ft} & \textbf{Available} & \textbf{ft} & \textbf{Margin (ft)} & \textbf{Result}\\\midrule
Takeoff Distance (TOD) & \num{2005} & TODA & \num{10302} & \num{8297} & OK\\
Takeoff ground roll & \num{1485} & TORA & \num{10302} & \num{8817} & OK\\
Landing Distance (LD) & \num{1688} & LDA & \num{10302} & \num{8614} & OK\\
\bottomrule\end{tabular}\caption{Cirrus SR22T: runway check, KICT runway 01L.}\end{table}
